\cvsection{Expériences professionnelles}

\begin{cventries}
  \cventry
    {Stagiaire ingénieur de recherche en IA}
    {INRIA}
    {Bordeaux, France}
    {Mars 2026 -- Septembre 2026}
    {
      \begin{cvitems}
        \item {Conception et développement en autonomie, avec le suivi de mes encadrants, d’une application web répondant au besoin d’aider les chercheurs à trouver des articles scientifiques.}
        \item {Développement du backend en Python/FastAPI, d’API REST documentées avec OpenAPI et de la persistance MongoDB.}
        \item {Développement d’une interface React/Next.js permettant aux chercheurs d’utiliser l’application et ses différents outils.}
        \item {Conteneurisation de composants avec Docker et mise en place d’un pipeline CI/CD sur GitHub pour automatiser leur intégration.}
        \item {Création d’un outil Python de collecte automatisée de données et de webhooks transmettant les résultats de tâches cron.}
        \item {Prise en charge d’une application web conteneurisée utilisant des modèles installés localement sur un Mac Studio, avec installation et configuration des modèles.}
        \item {Évaluation d’OpenClaw dans un environnement Docker isolé avec une attention portée aux risques liés à l’accès autonome aux outils et aux fichiers.}
        \item {Utilisation de Codex et de DeepSeek dans des workflows de développement assisté par l’IA pour accompagner le prototypage et l’amélioration du code.}
        \item {Définition d’un protocole expérimental, analyse des résultats et contribution à un article scientifique comparant plusieurs modèles.}
      \end{cvitems}
    }

  \cventry
    {Stagiaire de recherche en apprentissage par renforcement}
    {Université Karunya}
    {Coimbatore, Tamil Nadu, Inde}
    {Juin 2025 -- Septembre 2025}
    {
      \begin{cvitems}
        \item {Développement d’un modèle d’apprentissage par renforcement et de son environnement d’entraînement pour la navigation autonome d’un robot hospitalier.}
        \item {Déploiement et optimisation du logiciel et du modèle sur un Raspberry Pi embarqué sous Linux.}
      \end{cvitems}
    }
\end{cventries}
